% 3+1 Primzahlfamilien mod 210 — aus Rohtext in LaTeX gesetzt
\documentclass[11pt,a4paper]{article}
\usepackage[utf8]{inputenc}
\usepackage[T1]{fontenc}
\usepackage[ngerman]{babel}
\usepackage{amsmath, amssymb, amsthm}
\usepackage{booktabs}
\usepackage{microtype}
\usepackage[unicode=true]{hyperref}

\newtheorem*{observation}{Theorem / Observation}

\title{Primzahlvierlinge: die $3{+}1$-Struktur modulo $210$}
\author{}
\date{}

\begin{document}
\maketitle

\begin{center}
\fbox{\parbox{0.92\linewidth}{\centering
\textbf{Kernaussage.} Primzahlvierlinge erscheinen modulo $210$ in drei Hauptfamilien plus einer Ankerklasse.}}
\end{center}

\section{Hauptergebnis: die $3{+}1$-Struktur modulo $210$}

Für Primzahlvierlinge
\[
  (p,\, p+2,\, p+6,\, p+8)
\]
zeigt die Analyse modulo
\[
  210 = 2 \cdot 3 \cdot 5 \cdot 7
\]
eine sehr starke Einschränkung der möglichen Startwerte $p$. Von den $48$ invertierbaren Restklassen modulo $210$ bleiben für die in der Praxis beobachteten Starts effektiv nur vier zulässige Klassen:
\[
  \boxed{p \equiv 5,\, 11,\, 101,\, 191 \pmod{210}.}
\]

Davon ist
\[
  \boxed{p \equiv 5 \pmod{210}}
\]
die \textbf{Ankerfamilie}: sie enthält den kleinsten und in diesem Sinne ausgezeichneten Vierling
\[
  (5,\, 7,\, 11,\, 13).
\]
Die übrigen, statistisch tragenden Vierlinge verteilen sich auf die drei Hauptfamilien
\[
  \boxed{11,\, 101,\, 191 \pmod{210}.}
\]

\subsection*{Empirische Zahlen (Beispieldatensatz)}

\begin{table}[htbp]
  \centering
  \begin{tabular}{@{}rr@{}}
    \toprule
    Familie $p \bmod 210$ & Anzahl \\
    \midrule
    $5$   & $1$ \\
    $11$  & $352$ \\
    $101$ & $327$ \\
    $191$ & $320$ \\
    \bottomrule
  \end{tabular}
\end{table}

Damit lässt sich die Struktur als
\[
  \boxed{\{5\} \;\cup\; \{11, 101, 191\}}
\]
lesen, also
\[
  \boxed{\text{eine Ankerklasse plus drei statistische Familien.}}
\]

\section{Warum genau diese vier Klassen?}

Der Vierling $(p,\, p{+}2,\, p{+}6,\, p{+}8)$ darf modulo $2$, $3$, $5$ und $7$ keine der vier Zahlen in eine durch diese Primzahl teilbare Restklasse schicken. Modulo $210$ werden daher alle Startklassen ausgeschlossen, für die mindestens eine der Zahlen
\[
  p,\quad p+2,\quad p+6,\quad p+8
\]
durch $2$, $3$, $5$ oder $7$ teilbar wäre.

Übrig bleiben genau die vier oben genannten Klassen. Die Klasse $5$ ist dabei singulär (Anfangsvierling); die drei großen Familien tragen die wiederkehrende empirische Masse.

\section{Formulierung für eine englischsprachige Kurzfassung}

\begin{observation}
Let $Q(p) = (p, p+2, p+6, p+8)$ be a prime quadruplet. Then the starting prime $p$ belongs modulo $210$ to one of the four residue classes
\[
  p \equiv 5,\, 11,\, 101,\, 191 \pmod{210}.
\]
The residue class $5 \pmod{210}$ contains the exceptional initial quadruplet $Q(5) = (5,7,11,13)$. Apart from this anchoring class, the nontrivial population of prime quadruplets is distributed among the three families
\[
  \boxed{11,\, 101,\, 191 \pmod{210}.}
\]
Empirically, in the tested range, these three families are approximately equidistributed.
\end{observation}

\section{Interpretation im BM/EABC-Stil}

Die mod-$12$-Diskussion liefert das grobe EABC-Muster ($\mathrm{ABCE}$ oder $\mathrm{CEAB}$); mod $36$ eine dreifache Phasenfeinstruktur. Modulo $210$ ist die Einschränkung stärker:
\[
  \boxed{\text{Mod } 210 \text{ ist die erste Ebene, auf der die Vierlinge in drei echte statistische Familien zerfallen.}}
\]
Die Ankerklasse $5$ ist der Ursprungspunkt:
\[
  5 \longrightarrow (5,\, 7,\, 11,\, 13),
\]
während $11$, $101$, $191$ die tragende Wiederholstruktur bilden.

\section{Vorschlag für einen Abstract}

\begin{quote}
\textbf{Abstract.}
We study prime quadruplets $Q(p) = (p, p+2, p+6, p+8)$ through their residue structure modulo $210 = 2\cdot 3\cdot 5\cdot 7$. The main observation is that admissibility reduces the starting primes $p$ to exactly four residue classes,
\[
  p \equiv 5,\, 11,\, 101,\, 191 \pmod{210}.
\]
The class $p \equiv 5 \pmod{210}$ acts as an anchoring class containing the initial quadruplet $(5,7,11,13)$, while all subsequent nontrivial quadruplets fall into three statistical families $11,\, 101,\, 191 \pmod{210}$. In the tested range, these three families are approximately equidistributed. Additional tests on transition matrices, gap sequences, and center smoothness show no evidence for dynamical correlations beyond the local congruence constraints.
\end{quote}

\section{Leitsatz}

\begin{center}
\fbox{\parbox{0.92\linewidth}{
Primzahlvierlinge besitzen modulo $210$ eine $3{+}1$-Familienstruktur:\quad
\textit{eine Ankerfamilie $5$ und drei Hauptfamilien $11$, $101$, $191$.}}}
\end{center}

\end{document}
